%% ===========================================================================
%% REVISION INSERTIONS: blank-image control (Reviewers 9 and 11), all in BLUE.
%%
%% Paste the three blocks below into the v5 main.tex at the marked places.
%% Nothing else in the paper is changed. No new packages or macros are needed
%% (xcolor, booktabs, graphicx and the \qwen/\medvlm/\gemma/\medgemma/\bn macros
%% are already in the v5 preamble). Labels used: sec_blank_method, sec_res_blank,
%% tab_blank, tab_blank_groups. References used: sec_stats, sec_res_bangla.
%%
%% All numbers are recomputed from the saved per-item result CSVs
%% (BUS_BRA/Results/Blank Control, RadImageNet_VQA/Blank_Control,
%%  BanglaMedVQA/Blank_Control and the matching real-image result files).
%% ===========================================================================


%% ###########################################################################
%% BLOCK 1 -- METHODOLOGY
%% Place: after the subsection "Inference and answer parsing" (label
%% sec_inference) and before "\subsection{Metrics and statistics}".
%% ###########################################################################
\subsection{\textcolor{blue}{Blank-image control}}
\label{sec_blank_method}
{\color{blue}
To measure how much each answer depends on the image, we repeated all three tasks with the image replaced by a constant blank image and everything else unchanged: the same four models, items, prompts, decoding, answer parsing and metrics, and no retraining or tuning. For RadImageNet-VQA and BanglaMedVQA the input was a $512\times512$ solid white RGB image, passed in place of the real image for each of the 500 questions. For BUS-BRA we used both a white and a black $512\times512$ image. Because the input is identical for all 400 BUS-BRA items and decoding is deterministic, we ran five identical inferences per colour and model, observed the same output every time, and assigned that output to all 400 items; the BUS-BRA blank-image scores are therefore derived from five inferences per condition, not from 400 independent ones.

We compare each model's real-image and blank-image answers on the same items. Accuracy differences use the exact McNemar test and LAVE differences the paired permutation test of Section~\ref{sec_stats}, with a 95\% paired bootstrap interval. To ask whether a medical model uses the image more than its base model, we also compare their real-minus-blank LAVE differences item by item (bootstrap and permutation test as above). These checks are descriptive and are not among the six main hypotheses, so they are not Holm-adjusted. We did not run a shuffled-image control (real images paired with the wrong questions).
}


%% ###########################################################################
%% BLOCK 2 -- RESULTS
%% Place: at the end of the BanglaMedVQA results subsection (after the
%% paragraph ending "...against our 31.5\% on a different sample.") and
%% before "\subsection{Discussion}".
%% ###########################################################################
\subsection{\textcolor{blue}{Blank-image control: how much do the answers depend on the image?}}
\label{sec_res_blank}
{\color{blue}
Table~\ref{tab_blank} compares each model's score on the real images with its score on the blank image. On RadImageNet-VQA and BanglaMedVQA every model scored much lower with the blank image (RadImageNet-VQA accuracy: McNemar $p<10^{-10}$ for all four models; BanglaMedVQA LAVE: paired permutation $p<0.001$ for all four). On BUS-BRA the blank image changed the accuracy only for \medgemma{}.
}

\begin{table}[tb]
\color{blue}
\centering
\caption{Real-image $\rightarrow$ blank-image scores. Accuracy (\%) on BUS-BRA (chance 50\%) and RadImageNet-VQA (chance 40\% over the 500 questions), and mean LAVE (\%) on BanglaMedVQA. The blank image is white for RadImageNet-VQA and BanglaMedVQA; on BUS-BRA the white and black images gave identical results.}
\label{tab_blank}
\footnotesize
\setlength{\tabcolsep}{3pt}
\resizebox{\columnwidth}{!}{%
\begin{tabular}{@{}lccc@{}}
\toprule
\textbf{Model} & \textbf{BUS-BRA acc.} & \textbf{RadImageNet acc.} & \textbf{Bangla LAVE} \\
\midrule
\qwen{}-2B & 50.0 $\rightarrow$ 50.0 & 64.2 $\rightarrow$ 41.6 & 15.1 $\rightarrow$ 2.7 \\
\medvlm{} & 50.0 $\rightarrow$ 50.0 & 65.0 $\rightarrow$ 43.8 & 16.0 $\rightarrow$ 2.6 \\
\gemma{}~4B & 50.3 $\rightarrow$ 50.0 & 63.0 $\rightarrow$ 38.8 & 24.1 $\rightarrow$ 11.7 \\
\medgemma{}~4B & 57.5 $\rightarrow$ 50.0 & 61.2 $\rightarrow$ 38.4 & 31.5 $\rightarrow$ 7.2 \\
\bottomrule
\end{tabular}}
\end{table}

{\color{blue}
\textbf{RadImageNet-VQA: anatomy depends on the image, disease barely does.} With the blank image, overall accuracy fell to 38.4--43.8\%, and the 95\% interval of every model includes the 40\% chance level. The loss was concentrated in the anatomy groups (Table~\ref{tab_blank_groups}): anatomy MCQ fell from 80--88\% to 11--36\% and anatomy yes/no from 79--88\% to 51--57\%, a gain from seeing the image of 28 to 73 points. Disease questions were already close to chance with the real image and moved little: pathology MCQ was 24--25\% (chance 25\%) against 32--37\%, and abnormality and pathology yes/no were 50--54\% against 52--61\%. The gain from the image on disease questions was between $-2$ and $+13$ points for every model and group. Two blank-image anatomy MCQ scores lay below the 25\% chance level (\gemma{} 11\%, \medgemma{} 15\%) because these models favoured letter B (63 and 77 of 100 answers), which is correct for only 22 of the 100 questions.
}

\begin{table}[tb]
\color{blue}
\centering
\caption{RadImageNet-VQA accuracy (\%) by group, real $\rightarrow$ blank image (100 questions per group; chance 25\% for MCQ and 50\% for yes/no).}
\label{tab_blank_groups}
\footnotesize
\setlength{\tabcolsep}{3pt}
\resizebox{\columnwidth}{!}{%
\begin{tabular}{@{}lcccc@{}}
\toprule
 & \textbf{Qwen2} & \textbf{MedVLM} & \textbf{Gemma3} & \textbf{MedGemma} \\
\midrule
Anatomy, MCQ & 87 $\rightarrow$ 28 & 88 $\rightarrow$ 36 & 84 $\rightarrow$ 11 & 80 $\rightarrow$ 15 \\
Anatomy, yes/no & 88 $\rightarrow$ 55 & 86 $\rightarrow$ 57 & 81 $\rightarrow$ 52 & 79 $\rightarrow$ 51 \\
Abnormality, yes/no & 55 $\rightarrow$ 50 & 61 $\rightarrow$ 50 & 52 $\rightarrow$ 54 & 55 $\rightarrow$ 50 \\
Pathology, yes/no & 57 $\rightarrow$ 51 & 57 $\rightarrow$ 51 & 61 $\rightarrow$ 53 & 60 $\rightarrow$ 51 \\
Pathology, MCQ & 34 $\rightarrow$ 24 & 33 $\rightarrow$ 25 & 37 $\rightarrow$ 24 & 32 $\rightarrow$ 25 \\
\bottomrule
\end{tabular}}
\end{table}

{\color{blue}
\textbf{Which answer preferences come from the image?} The control separates preferences that depend on the input from those that do not. The option-letter preference does not: on pathology MCQ the most chosen letter with the blank image was the same as with the real image for all four models (C for the Qwen pair, B for the Gemma pair), and \gemma{} chose B 54 times instead of 35, so this preference behaves like a language prior inherited from the base model. The yes/no preferences do depend on the image. The share of ``yes'' answers over the 300 yes/no questions fell from 54\%, 53\%, 74\% and 42\% to 12\%, 13\%, 30\% and 14\% (\qwen{}, \medvlm{}, \gemma{}, \medgemma{}), and on abnormality questions no model said ``yes'' more than 4 times in 100 (real image: 77, 69, 98 and 39). \gemma{}'s habit of saying ``yes'' to almost every abnormality question therefore appears only when a scan is shown; with a blank image all four models mostly answered ``no''.

BUS-BRA shows the same pattern. With a blank image, white or black and in every repetition, all four models answered ``benign'' for all 400 cases. \qwen{} and \medvlm{} gave this answer for the real images as well (identical on all 400 cases), so their BUS-BRA output does not depend on the image at all. \gemma{} answered ``malignant'' for 391 real images but ``benign'' for the blank one (2\% identical answers), so its ``malignant'' preference is triggered by the ultrasound image rather than by the prompt. \medgemma{} gave the same answer as with the blank image on 66\% of cases, and its 57.5\% accuracy on real images exceeds the 50.0\% of the blank image (82 cases correct only with the real image, 52 only with the blank one; exact McNemar $p=0.012$, unadjusted).

\textbf{BanglaMedVQA: \medgemma{} depends on the image more than \gemma{}.} Mean LAVE fell by 12.4, 13.4, 12.4 and 24.3 points for \qwen{}, \medvlm{}, \gemma{} and \medgemma{} (95\% bootstrap intervals 9.7--15.3, 10.8--16.2, 9.0--15.8 and 20.7--28.0), and exact match was 0 of 500 for every model except \gemma{} (2 of 500). The blank-image answers were nearly constant within a category. \gemma{} answered ``X-ray'' to 94 of the 100 modality questions, \medgemma{} answered ``CT'' to all 100 and ``brain'' to 98 of the 100 organ questions, and \qwen{} and \medvlm{} named no modality at all (0 of 100 correct by the rule-based check of Section~\ref{sec_res_bangla}; typical answers ``this is a medical image'' or ``digital''); \qwen{} answered \bn{লাইন} (``line'') to 33 of the 100 organ questions. A constant answer earns the share of references it happens to match: X-ray is correct for 56\% of the modality references, so \gemma{} kept a modality LAVE of 48.7 on the blank image (real image 60.7), whereas \medgemma{} fell from 73.3 to 7.3 and \qwen{} and \medvlm{} from 55.8 and 56.5 to 1.6 and 2.3. Organ LAVE fell from 47.1 to 0.5 for \gemma{} and from 63.1 to 21.4 for \medgemma{}. On the three specialised categories (abnormality, condition, position) LAVE was 0.0--11.1 with the real images and 0.3--7.3 with the blank image. \gemma{} also returned an empty answer for 125 of the 500 blank-image questions (11 with real images), which lowers its blank-image scores.

The real-minus-blank LAVE of \medgemma{} is 11.9 points larger than that of \gemma{} (95\% CI 7.2--16.5; paired permutation $p<0.001$), whereas \medvlm{} and \qwen{} differ by only 0.9 points ($-0.6$ to $2.4$; $p=0.23$). With the blank image \medgemma{} scored lower than \gemma{} (7.2 against 11.7), so its 7.4-point advantage with real images is not a language prior: it appears only when an image is shown.

\textbf{What the control changes.} It leaves the paired comparison unchanged: \medvlm{} remains indistinguishable from \qwen{}, and the \medgemma{} advantage in Bangla is image-dependent. It sharpens two statements. Anatomy and modality answers depend on the image, whereas answers to disease questions were only slightly better with the real image than with a blank one. And the answer preferences are of two kinds: the option-letter preference is independent of the image, while the ``yes'' and ``malignant'' preferences appear only when a scan is shown. A blank image is an unnatural input and we used one white image per task (white and black on BUS-BRA), so the control does not replace a shuffled-image test.
}


%% ###########################################################################
%% BLOCK 3 -- CONCLUSIONS (one sentence; needed for consistency)
%% In the Conclusions, the last sentence currently ends
%%   "Future work should add per-model prompt tuning, blank-image controls,
%%    English versions of the Bangla questions and larger models."
%% Replace that sentence with the following (the old one would now be false,
%% since the blank-image controls have been run). Paste it inline, with no blank
%% line before it, so the Conclusions stay one paragraph:
%% ###########################################################################
{\color{blue}
Blank-image controls showed that the models' anatomy and modality answers depend on the image, whereas answers to disease questions were only slightly better with the real image than with a blank one. Future work should add per-model prompt tuning, shuffled-image controls, English versions of the Bangla questions and larger models.
}
